% MANUSCRIPT.tex, the manuscript skeleton for the asymmetric capital buffer
% paper. Created 10 October 2026 as pass 1 of the plan in TODO.md (item PLAN).
%
% PURPOSE
%   The manuscript. Four tables carry the body (introduction, model,
%   literature, conclusions), and the appendices carry every proof.
%   00_document/PROOFS_v2.tex is retired into this file at pass 6, with
%   RETIREMENT.md recording where each part of it went.
%
% STRUCTURE THE CHECKER RELIES ON (checks/verify_manuscript.py)
%   Rows are written only through the four row macros below, and proofs only
%   through the mproof environment. IDs are K1, K2, ... (introduction),
%   M1, ... (model), L1, ... (literature) and C1, ... (conclusions).
%     \krow{ID}{headline}{model rows}{scope note}
%     \mrow{ID}{claim in words}{statement in symbols}{Lean theorems}
%     \lrow{ID}{bib key}{what we rely on}{verbatim quote}{page}{model rows}
%     \crow{ID}{statement}{model or literature rows}
%     \begin{mproof}{model row}{Lean theorems} ... \end{mproof}
%   Lists of rows or theorems are comma-separated. Lean theorem names are the
%   fully qualified names declared under lean/mathlib/.
%
% RULES THE CHECKER ENFORCES
%   Introduction rows point to model rows only. Conclusion rows point to model
%   or literature rows only. Every model row names Lean theorems that exist in
%   an audited Lean file and has exactly one proof in Appendix A naming the
%   same theorems. Every literature row's quote appears verbatim in the
%   paper's lit/ record. In every proof, plain English is at most 30% of the
%   text, measured as prose words over prose words plus mathematical tokens.

\documentclass[11pt]{article}
\usepackage[T1]{fontenc}
\usepackage[utf8]{inputenc}
\usepackage{mathptmx}
\usepackage{microtype}
\usepackage[a4paper,margin=1.8cm]{geometry}
\usepackage{amsmath,amssymb,amsthm}
\usepackage{booktabs}
\usepackage{array}
\usepackage{longtable}
\usepackage{url}
\usepackage[hidelinks]{hyperref}
\usepackage{natbib}

\emergencystretch=1.5em

\newcolumntype{P}[1]{>{\raggedright\arraybackslash}p{#1}}

\urlstyle{tt}
\makeatletter
\g@addto@macro\UrlBreaks{\do\.\do\_\do\,}
\makeatother
\newcommand{\lean}[1]{\path{#1}}
\newcommand{\rowref}[1]{#1}
\newcommand{\unv}[1]{\textit{Unverified:} #1}
\newcommand{\pend}[1]{\textit{Pending:} #1}

\newcommand{\krow}[4]{#1 & #2 & \rowref{#3} & #4 \\ \addlinespace}
\newcommand{\mrow}[4]{#1 & #2 & #3 & \lean{#4} & p.~\pageref{proof:#1} \\ \addlinespace}
\newcommand{\lrow}[6]{#1 & \citet{#2} & #3 & ``#4'' & #5 & \rowref{#6} \\ \addlinespace}
\newcommand{\crow}[3]{#1 & #2 & \rowref{#3} \\ \addlinespace}

\newenvironment{mproof}[2]
  {\subsection*{Proof of #1}\label{proof:#1}%
   {\raggedright\noindent\textit{Lean:} \lean{#2}.\par}\smallskip}
  {\par\hfill$\square$\par\medskip}

\title{Asymmetric Capital Buffer Control}
\author{}
\date{Manuscript skeleton, 10 October 2026}

\begin{document}
\maketitle

\section{Introduction}

The question is how an asymmetric penalty on a capital shortfall shapes a
bank's recapitalisation policy and the volatility of its capital buffer.
The model treats recapitalisation as impulse control under a fixed issuance
cost and a regulatory cap on risky exposure.

\footnotesize
\begin{longtable}{@{}P{0.8cm}P{7.0cm}P{1.9cm}P{6.4cm}@{}}
\caption{Headlines.}\label{tab:intro}\\
\toprule
ID & Headline & Model rows & Scope note \\
\midrule
\endfirsthead
\toprule
ID & Headline & Model rows & Scope note \\
\midrule
\endhead
\bottomrule
\end{longtable}
\normalsize

\section{Model}

\footnotesize
\begin{longtable}{@{}P{0.8cm}P{3.8cm}P{6.3cm}P{3.7cm}P{1.1cm}@{}}
\caption{Model results.}\label{tab:model}\\
\toprule
ID & Claim & Statement & Lean & Proof \\
\midrule
\endfirsthead
\toprule
ID & Claim & Statement & Lean & Proof \\
\midrule
\endhead
\bottomrule
\end{longtable}
\normalsize

\section{Literature}

\footnotesize
\begin{longtable}{@{}P{0.8cm}P{2.2cm}P{3.5cm}P{5.5cm}P{1.3cm}P{1.9cm}@{}}
\caption{Papers read.}\label{tab:lit}\\
\toprule
ID & Paper & What we rely on & Quote & Page & Rows \\
\midrule
\endfirsthead
\toprule
ID & Paper & What we rely on & Quote & Page & Rows \\
\midrule
\endhead
\bottomrule
\end{longtable}
\normalsize

\section{Conclusions}

Table~\ref{tab:concl} holds the conclusions that rest on model rows and on
papers read in full. Table~\ref{tab:pending} holds every claim that a result
is new and every other claim still waiting, and each of its rows names what
it waits on.

\footnotesize
\begin{longtable}{@{}P{0.8cm}P{12.3cm}P{3.4cm}@{}}
\caption{Conclusions.}\label{tab:concl}\\
\toprule
ID & Statement & Rows \\
\midrule
\endfirsthead
\toprule
ID & Statement & Rows \\
\midrule
\endhead
\bottomrule
\end{longtable}

\begin{longtable}{@{}P{0.8cm}P{12.3cm}P{3.4cm}@{}}
\caption{Pending conclusions.}\label{tab:pending}\\
\toprule
ID & Statement & Rows \\
\midrule
\endfirsthead
\toprule
ID & Statement & Rows \\
\midrule
\endhead
\bottomrule
\end{longtable}
\normalsize

\appendix

\section{Proofs}
\label{app:proofs}

\paragraph{Notation.}

\section{Evidence and verification}
\label{app:evidence}

\paragraph{Toolchain.} Lean \path{v4.32.0-rc1} with Mathlib at tag
\path{v4.32.0-rc1}, commit \path{360da6fa66c1}, pinned in \path{lean/mathlib/lake-manifest.json}. Every file under
\path{lean/mathlib/} is a build root. A theorem passes the axiom audit when it
depends on no axiom outside \path{propext}, \path{Classical.choice} and
\path{Quot.sound}.

\paragraph{Checks.} \path{verify.sh} runs \path{checks/verify_manuscript.py},
\path{checks/verify_mathlib.py} and \path{checks/verify_docs.py}, and
writes their output to \path{VERIFICATION.md}.
\path{lit/verify_lit.sh} runs the suite of every literature record.

\paragraph{What Lean does not carry.}

\paragraph{Illustrations.} No row rests on SymPy or on numerical sampling.

\section{Terminology}
\label{app:terms}

\footnotesize
\begin{longtable}{@{}P{3.4cm}P{4.2cm}P{8.6cm}@{}}
\toprule
Term & Symbol & Meaning \\
\midrule
\endfirsthead
\toprule
Term & Symbol & Meaning \\
\midrule
\endhead
\bottomrule
\end{longtable}
\normalsize

\section{Primitives and scope}
\label{app:primitives}

\section{Refuted conjectures}
\label{app:refuted}

\end{document}
